\noindent\textbf{Background.}
Machine learning models for chronic kidney disease risk prediction regularly achieve strong discrimination on internal test sets. Calibration assessment and uncertainty quantification are far less common, leaving clinicians without reliable information about whether probability outputs are trustworthy. No published study has jointly evaluated all three dimensions (calibration, uncertainty, and structured deployment readiness) on a common model suite with external clinical validation.

\medskip
\noindent\textbf{Objective.}
To evaluate five classifiers across calibration quality, conformal prediction coverage, and an eight-criterion deployment readiness framework on both internal and external data.

\medskip
\noindent\textbf{Methods.}
Five classifiers (logistic regression, random forest, XGBoost, support vector machine with Platt scaling, Gaussian naive Bayes) were trained on the UCI CKD dataset (400 patients, 62.5\% CKD). A distributional stress-test used the open-access MIMIC-IV demo cohort (97 patients, 23.7\% CKD) to evaluate model behaviour under prevalence shift and feature missingness. Calibration was assessed before and after Platt scaling and isotonic regression, quantified by Expected Calibration Error and Brier Score. Predictive uncertainty was measured through split conformal prediction targeting 90\% marginal coverage. An eight-criterion deployment readiness framework evaluated discrimination, calibration stability, coverage transfer, subgroup equity, and reproducibility.

\medskip
\noindent\textbf{Results.}
All five models achieved AUROC 1.00 on the UCI test set. Post-isotonic ECE fell to 0.000--0.022 internally. On MIMIC-IV, AUROC dropped to 0.48--0.58, ECE rose to 0.68--0.76, and conformal coverage collapsed from 0.80--0.98 (UCI) to 0.21--0.25, well below the 90\% target. No model passed the deployment checklist; scores ranged from 2 to 4 out of 16.

\medskip
\noindent\textbf{Conclusion.}
Near-perfect internal performance did not survive distributional shift. Calibration stability and conformal coverage transfer should be evaluated before any clinical ML model moves toward deployment, even when internal metrics appear strong.
